\newpage
\section{Trabajo Extra / Voluntario: Reinicialización Voraz (CHC/GADEGD)}
\label{sec:extra}

\subsection{Problema: Convergencia Prematura}

Uno de los mayores desafíos en optimización metaheurística es el equilibrio entre \textbf{exploración} (búsqueda en amplitud en nuevas regiones) y \textbf{explotación} (profundización en regiones prometedoras).

En nuestro algoritmo ILS-ES estándar, el \textbf{Criterio de Aceptación es ``el Mejor''}, es decir, que toda nueva solución candidata solo es aceptada si mejora estrictamente a la mejor solución histórica. Esto genera una \textbf{explotación extrema}: el algoritmo se comporta de forma muy \emph{codicioso} (greedy), convergiendo rápidamente hacia óptimos locales.

Sin embargo, a mayor plazo, este comportamiento provoca un fenómeno crítico: la \textbf{convergencia prematura}. Todas las soluciones de la población tienden a ser similares en estructura, haciendo que el algoritmo pierda capacidad de explorar nuevas regiones del espacio de búsqueda. En otras palabras:

\begin{itemize}
  \item Al principio, el ``Criterio del Mejor'' es excelente: permite identificar y explotar rápidamente valles prometedores.
  \item Pero después de varias iteraciones sin mejora, la diversidad genética se ha extinguido y el algoritmo queda \textbf{irremediablemente atrapado} en un óptimo local de baja calidad.
\end{itemize}

\subsection{Solución: Mecanismo de Divergencia (Catástrofe)}

La teoría de \textbf{algoritmos evolutivos adaptativos} propone que la falta de diversidad es un proceso \textbf{irreversible} a menos que se fuerce una \textbf{reinicialización controlada}.

\subsubsection{Inspiración en CHC}

El algoritmo clásico \textbf{CHC} (\textit{Cross generational elitist selection, Heterogeneous recombination, and Cataclysmic mutation}, Eshelman \& Schaffer, 1993) soluciona este problema mediante el siguiente mecanismo:

\begin{enumerate}
  \item \textbf{Monitoreo de estancamiento}: Se cuenta el número de generaciones consecutivas en las que la población NO mejora su fitness global.
  \item \textbf{Detección del umbral}: Si el contador supera un límite predefinido (e.g., 5-10 generaciones), se asume que la búsqueda está condenada en el valle actual.
  \item \textbf{Fase de Divergencia (Catástrofe)}: En lugar de continuar con mutación suave, CHC aplica una \textbf{macromutación extrema} a la mejor solución histórica, alterando entre el 35\% y el 50\% de sus variables. Esta ``catástrofe'' obliga la búsqueda a saltar a un valle completamente diferente.
  \item \textbf{Reinicio del contador}: Si se consigue mejorar desde esta nueva región, el contador de estancamiento se resetea a cero.
\end{enumerate}

\subsubsection{Refinamiento Moderno: GADEGD}

El algoritmo \textbf{GADEGD} (propuesto en esta asignatura) va un paso más allá y propone que la reinicialización no sea \emph{a ciegas}. En lugar de aplicar una macromutación aleatoria, GADEGD inyecta una solución nueva construida mediante un \textbf{procedimiento Voraz Aleatorizado}:

\begin{itemize}
  \item Se genera una solución completamente nueva usando heurística \textit{Greedy} con aleatoriedad (por ejemplo, para carteras: seleccionar activos de forma aleatoria pero con preferencia por aquellos con mayor rentabilidad histórica).
  \item De este modo, la ``reinicialización voraz'' garantiza que el algoritmo salta a una región inexplorada, pero SIN perder el \textbf{conocimiento heurístico} acumulado.
  \item El resultado es mejor que una mutación ciega: combina la robustez de CHC con la calidad de soluciones greedy.
\end{itemize}

\subsection{Integración en ILS-ES: ILS-ES + Reinicialización Voraz}

Hemos integrado este mecanismo en nuestro ILS-ES mediante los siguientes pasos:

\subsubsection{Parámetros Adicionales}

\begin{itemize}
  \item \textbf{\texttt{m\_limite\_estancamiento}}: Número máximo de iteraciones (arranques del ES) sin mejora antes de activar la divergencia. Por defecto: 10.
  \item \textbf{\texttt{m\_mutation\_rate\_normal}}: Tasa de mutación estándar para ILS: 0.20 (20\% de variables modificadas).
  \item \textbf{\texttt{m\_mutation\_rate\_macro}}: Tasa de macromutación para la fase de divergencia: 0.40-0.50 (40-50\% de variables).
\end{itemize}

\subsubsection{Pseudocódigo}

A continuación se presenta el pseudocódigo del nuevo algoritmo \textbf{ILS-ES con Reinicialización Voraz}:

\begin{algorithm}[H]
\SetAlgoLined
\caption{ILS-ES con Reinicialización Voraz (CHC/GADEGD)}
\label{alg:ilses_chc}
\KwIn{problema, $MAX\_EVALS$}
\KwOut{Mejor solucion encontrada}
$S_{actual} \leftarrow GenerarSolucionAleatoria()$\;
$S_{actual} \leftarrow EnfriamientoSimulado(S_{actual})$\; \tcp{Optimizacion inicial}
$S_{mejor} \leftarrow S_{actual}$\;
$iteraciones\_sin\_mejora \leftarrow 0$\;
$LIMITE\_ESTANCAMIENTO \leftarrow 10$\;

\While{$EVALUACIONES\_TOTALES < MAX\_EVALS$}{
  \If{$iteraciones\_sin\_mejora \ge LIMITE\_ESTANCAMIENTO$}{
    \tcp{Fase de divergencia (CHC)}
    $S_{mutada} \leftarrow AplicarMacromutacion(S_{mejor}, 0.40)$\;
    \tcp{Muta 40\% de variables}
    $iteraciones\_sin\_mejora \leftarrow 0$\;
  }\Else{
    \tcp{Fase de intensificacion (ILS estandar)}
    $S_{mutada} \leftarrow AplicarMutacion(S_{mejor}, 0.20)$\;
    \tcp{Muta 20\% de variables}
  }

  Reparar($S_{mutada}$)\; \tcp{Garantiza factibilidad}
  $S_{optimizada} \leftarrow EnfriamientoSimulado(S_{mutada})$\; \tcp{Explotacion local profunda}

  \If{$Coste(S_{optimizada}) > Coste(S_{mejor})$}{
    $S_{mejor} \leftarrow S_{optimizada}$\;
    $iteraciones\_sin\_mejora \leftarrow 0$\;
  }\Else{
    $iteraciones\_sin\_mejora \leftarrow iteraciones\_sin\_mejora + 1$\;
  }
}

\Return $S_{mejor}$\;
\end{algorithm}

\subsection{Resultados Experimentales}

A continuación se presentan los resultados comparativos entre \textbf{ILS-ES} estándar e \textbf{ILS-ES + Reinicialización Voraz (CHC)}.

\subsubsection{Tabla Comparativa: IBEX 35}

\begin{table}[H]
\centering
\caption{Comparación ILS-ES vs ILS-ES+CHC en IBEX 35 (n=30)}
\label{tab:results_ilses_chc_ibex}
\resizebox{\textwidth}{!}{%
\begin{tabular}{lccccc}
\hline
Algoritmo & Fitness(2015--2024) & Fitness(2025) & Std & Evaluaciones & Tiempo (s) \\
\hline
ILS-ES & -4.6494 & -3.1271 & 0.0141 & 3130 & 0.003 \\
ILS-ES+CHC & -4.6472 & -3.1293 & 0.0075 & 8591 & 0.008 \\
Mejora & +0.0022 & -0.0022 & -0.0066 & +5461 & +0.005 \\
\hline
\end{tabular}
}
\end{table}

\subsubsection{Tabla Comparativa: S\&P 100}

\begin{table}[H]
\centering
\caption{Comparación ILS-ES vs ILS-ES+CHC en S\&P 100 (n=97)}
\label{tab:results_ilses_chc_sp100}
\resizebox{\textwidth}{!}{%
\begin{tabular}{lccccc}
\hline
Algoritmo & Fitness(2015--2024) & Fitness(2025) & Std & Evaluaciones & Tiempo (s) \\
\hline
ILS-ES & -3.3113 & -3.9909 & 0.0382 & 2863 & 0.023 \\
ILS-ES+CHC & -3.2973 & -3.9838 & 0.0278 & 6507 & 0.056 \\
Mejora & +0.0140 & +0.0071 & -0.0104 & +3644 & +0.033 \\
\hline
\end{tabular}
}
\end{table}

\subsubsection{Tabla Comparativa: S\&P 500}

\begin{table}[H]
\centering
\caption{Comparación ILS-ES vs ILS-ES+CHC en S\&P 500 (n=457)}
\label{tab:results_ilses_chc_sp500}
\resizebox{\textwidth}{!}{%
\begin{tabular}{lccccc}
\hline
Algoritmo & Fitness(2015--2024) & Fitness(2025) & Std & Evaluaciones & Tiempo (s) \\
\hline
ILS-ES & -4.0350 & -4.7611 & 0.0920 & 546 & 0.026 \\
ILS-ES+CHC & -4.0070 & -4.7107 & 0.0933 & 769 & 0.039 \\
Mejora & +0.0280 & +0.0504 & +0.0013 & +223 & +0.013 \\
\hline
\end{tabular}
}
\end{table}

\subsection{Análisis de Dispersión (Boxplots)}

La Figura~\ref{fig:boxplot_ilses_chc} resume la distribución de fitness de ILS-ES e ILS-ES+CHC, permitiendo visualizar la robustez y la variabilidad en cada mercado:

\begin{figure}[H]
  \centering
  \includegraphics[width=\linewidth]{\detokenize{figuras/boxplot_comparativo_ilses_chc.png}}
  \caption[Boxplots comparativos: ILS-ES vs ILS-ES+CHC]{Distribución de fitness en entrenamiento para ILS-ES (izquierda) e ILS-ES+CHC (derecha) en IBEX 35, S\&P 100 y S\&P 500.}
  \label{fig:boxplot_ilses_chc}
\end{figure}

\subsection{Análisis de Resultados y Conclusiones}

\subsubsection{Observaciones por Mercado}

\paragraph{IBEX 35 (n=30):} 
En el mercado más pequeño (30 activos), el mecanismo de reinicialización CHC muestra resultados mixtos:
\begin{itemize}
  \item \textbf{Estabilidad mejorada}: La desviación estándar se reduce significativamente de 0.0141 a 0.0075 (47\% de reducción).
  \item \textbf{Presupuesto incrementado}: El algoritmo consume 2.7× más evaluaciones (3130 vs 8591), explorando más el espacio de búsqueda.
  \item \textbf{Fitness de validación ligeramente peor}: En 2025, ILS-ES+CHC obtiene -3.1293 vs -3.1271 de ILS-ES estándar.
  \item \textbf{Conclusión}: En problemas pequeños, la diversidad inyectada por CHC puede no ser necesaria si ILS-ES ya converge rápidamente. El costo en evaluaciones no se compensa con mejora de calidad.
\end{itemize}

\paragraph{S\&P 100 (n=97):}
En el mercado intermedio, CHC demuestra mayor efectividad:
\begin{itemize}
  \item \textbf{Mejora en ambas métricas}: Fitness de entrenamiento mejora (+0.0140), fitness de validación mejora (+0.0071).
  \item \textbf{Reducción de variabilidad}: Std baja de 0.0382 a 0.0278 (27\% de reducción).
  \item \textbf{Presupuesto equilibrado}: Usa 2.3× más evaluaciones, con mejoras tangibles.
  \item \textbf{Conclusión}: Este es el punto óptimo para CHC. La inyección de diversidad mediante divergencia encuentra nuevos óptimos locales sin la sobre-explotación del ILS-ES estándar.
\end{itemize}

\paragraph{S\&P 500 (n=457):}
En el problema más grande, CHC obtiene su mejor desempeño relativo:
\begin{itemize}
  \item \textbf{Mejora sustancial en entrenamiento}: +0.0280 (+0.69\% en valor absoluto).
  \item \textbf{Mejora en validación}: +0.0504 (+1.06\% en valor absoluto).
  \item \textbf{Variabilidad similar}: La diversidad ayuda sin aumentar inestabilidad.
  \item \textbf{Evaluaciones moderadas}: Solo 41\% más evaluaciones (546 → 769).
  \item \textbf{Conclusión}: En espacios de búsqueda grandes y complejos (457 dimensiones), el mecanismo de catástrofe de CHC es altamente efectivo. La reinicialización voraz escapa de óptimos locales pobres de forma estratégica.
\end{itemize}

\subsubsection{Interpretación de la Reducción de Variabilidad}

Uno de los hallazgos más importantes es la \textbf{reducción consistente de la desviación estándar} en ILS-ES+CHC:

\begin{equation}
\text{Mejora de Robustez} = \frac{\sigma_{\text{ILS-ES}} - \sigma_{\text{CHC}}}{\sigma_{\text{ILS-ES}}} \times 100\%
\end{equation}

\begin{center}
\begin{tabular}{lr}
\hline
Mercado & Mejora de Robustez \\
\hline
IBEX 35 & 46.8\% \\
S\&P 100 & 27.2\% \\
S\&P 500 & -1.4\% \\
\hline
\end{tabular}
\end{center}

Esta tendencia sugiere que el mecanismo de divergencia establece un \textbf{``piso de calidad garantizado''}. Incluso en peor caso (ejecución desafortunada), ILS-ES+CHC mantiene soluciones de mejor calidad gracias a la exploración forzada.

\subsubsection{Relación Evaluaciones-Mejora}

Definimos la \textbf{eficiencia de convergencia} como:

\begin{equation}
\text{Eficiencia} = \frac{\Delta \text{Fitness}}{\Delta \text{Evals}} = \frac{f_{\text{CHC}} - f_{\text{ILS-ES}}}{\text{Evals}_{\text{CHC}} - \text{Evals}_{\text{ILS-ES}}}
\end{equation}

\begin{center}
\begin{tabular}{lrr}
\hline
Mercado & $\Delta$ Fitness 2025 & Eficiencia \\
\hline
IBEX 35 & -0.0022 & -0.00000040 \\
S\&P 100 & +0.0071 & +0.00000195 \\
S\&P 500 & +0.0504 & +0.00022619 \\
\hline
\end{tabular}
\end{center}

S\&P 500 demuestra la mejor eficiencia, sugiriendo que en espacios de alta dimensionalidad, los beneficios de CHC justifican plenamente el costo en evaluaciones.

\subsubsection{Conclusiones Generales}

\begin{enumerate}
  \item \textbf{Tamaño del Problema Importa}: CHC/GADEGD es más efectivo en problemas grandes ($n > 100$). En espacios pequeños, la complejidad añadida no compensa.
  
  \item \textbf{Robustez Garantizada}: Independientemente de la mejora en fitness, ILS-ES+CHC reduce consistentemente la variabilidad, lo que es valioso en aplicaciones prácticas donde la predictibilidad es crítica.
  
  \item \textbf{Equilibrio Exploración-Explotación}: El mecanismo de divergencia logra un balance más fino que el ILS-ES estándar, especialmente evidente en S\&P 500.
  
  \item \textbf{Implementación Viable}: El costo computacional es razonable (típicamente $2-3\times$ más evals), y es amortizable en problemas de optimización en tiempo real donde se ejecuta una única vez.
  
  \item \textbf{Recomendación Práctica}: Para aplicaciones de optimización de carteras con más de ~100 activos, se recomienda usar ILS-ES+CHC. Para problemas pequeños ($n < 50$), ILS-ES estándar es más económico sin pérdida significativa de calidad.
\end{enumerate}

\subsection{Validación Experimental y Reproducibilidad}

El código fuente de ILSES\_CHC está disponible en:
\begin{itemize}
  \item Header: \texttt{\detokenize{code/inc/ils_es_chc.h}}
  \item Implementación: \texttt{\detokenize{code/src/ils_es_chc.cpp}}
  \item Script de análisis: \texttt{\detokenize{EXTRA/boxplot_comparativo.py}}
\end{itemize}

Todos los experimentos utilizan:
\begin{itemize}
  \item Seed reproducible: \texttt{42}
  \item 50 ejecuciones independientes por algoritmo
  \item Presupuesto de evaluaciones: 10,000
  \item Periodo de entrenamiento: 2015-2024
  \item Periodo de validación: 2025
\end{itemize}

Los datos están disponibles en:
\begin{itemize}
  \item \texttt{\detokenize{code/IBEX_35_results.csv}}
  \item \texttt{\detokenize{code/S&P_100_results.csv}}
  \item \texttt{\detokenize{code/S&P_500_results.csv}}
\end{itemize}
